% q0038.tex — Same Notes, Different Song
\question{
  id        = {0038},
  title     = {Same Notes, Different Song},
  source    = {OBECON},
  year      = {2026},
  round     = {Seletiva IEO -- Phase C},
  language  = {English},
  topic     = {Macro},
  subtopic  = {Unemployment / Composition Effects},
  type      = {objective},
  statement = {%
    Obeconomia has three age groups: Youth (50M), Adults (30M), Elderly (20M)
    with unemployment rates 10\%, 3\%, 7\% respectively, giving an overall
    unemployment rate of 7.3\%. After demographic shifts to Youth (20M),
    Adults (50M), Elderly (30M) with unchanged group-specific rates, what
    happens to overall unemployment?
  },
  choices   = {%
    \item Overall unemployment rises because within-group rates rise.
    \item Overall unemployment stays constant at 7.3\%.
    \item Overall unemployment falls from 7.3\% to 5.6\% due to composition
          effects.
    \item Overall unemployment increases from 5.6\% to 7.3\% due to
          composition effects.
  },
  answer    = {C},
  solution  = {%
    The new overall unemployment rate is:
    \[
      \frac{0.10 \times 20 + 0.03 \times 50 + 0.07 \times 30}{100}
      = \frac{2 + 1.5 + 2.1}{100} = 5.6\%.
    \]
    The shift toward the lower-unemployment adult group reduces aggregate
    unemployment despite unchanged group-specific rates --- a pure composition
    effect. (C) is correct.
  },
}
